% ch02 — Differential Geometry for GR
\chapter{Differential Geometry for General Relativity}
\label{ch:differential-geometry}

This chapter develops the differential-geometric machinery required to
formulate general relativity. We move from the abstract definition of a
manifold to the concrete tools---covariant derivatives, Christoffel symbols,
and curvature tensors---that appear directly in the codebase's metric and
geodesic modules.

\section{Manifolds, Tangent Spaces, and One-Forms}
\label{sec:dg-manifolds}

We define smooth manifolds, tangent vectors as directional derivatives, and
dual one-forms, establishing the arena on which tensor fields live.

\section{The Metric Tensor and Index Operations}
\label{sec:dg-metric}
\coderef{Math.Tensor}{MetricTensor}

We introduce the metric tensor as the object that defines distances, angles,
and the musical isomorphism between vectors and one-forms.

\section{Covariant Derivatives and Christoffel Symbols}
\label{sec:dg-covariant}
\coderef{Math.Tensor}{Christoffel}

We define the Levi-Civita connection, derive the Christoffel symbols from the
metric, and show how covariant differentiation generalises partial
differentiation to curved spaces.

\section{The Geodesic Equation}
\label{sec:dg-geodesic}

We derive the geodesic equation from parallel transport and from the
variational principle, connecting the geometry to the equations of motion
integrated by the ray tracer.

\section{Curvature Tensors}
\label{sec:dg-curvature}

We define the Riemann tensor via the commutator of covariant derivatives,
contract it to obtain the Ricci tensor and scalar, and state the symmetries
that reduce the number of independent components.

\paragraph{Key References.}
The presentation follows \cite{Carroll2004} and \cite{Wald1984}; notation
conventions are consistent with \cite{MTW1973}.
